% ! TeX root = ./main.tex

\section{Designing a Placement Policy}\label{sec:amr:design}

With the reliable and interpretable telemetry foundation established in \S\ref{sec:amr:tele}, we turned
to the final challenge: using this runtime information to design effective placement policies.
Validated measurements confirmed that MPI collective synchronization dominated runtime, accounting
for 35\%--50\% of total time and increasing with scale (512--4096 ranks), while direct
communication and redistribution overhead remained below 10\%.
This cost was clearly a downstream effect of compute time variability between blocks.
Mitigating the impact of this measured variability became the primary placement objective.

Guided by our codes---which, in the worst case, trigger refinement every five 250\,ms
timesteps---we set a 50\,ms placement computation budget to cap overhead at 5\% of the total
runtime.
This constraint ruled out complex optimization objectives such as maximizing communication-compute
overlap or encoding network topology, due to limited expected payoff and high constraint
complexity.
Instead, we focused on two key optimization dimensions:

\parahead{Compute Load
  Balance} Directly minimizing the variance in per-rank compute load on the basis of measured block
costs was the obvious objective to reduce straggler delays.
% Guided by our codes, which in the worst case refined every five 250\,ms timesteps, we set a target
% of 50\,ms as a guiding constraint to keep placement computation overhead under 5\% of the total
% runtime. Within this runtime constraint, we decided to explore two key dimensions: Achieving this
% required navigating a tradeoff between two key factors, under the constraint that each placement
% decision must complete in under 50\,ms:
This problem---formally known as \emph{makespan minization}---is
NP-hard~\cite{Garey1979:Intractability}, necessitating efficient heuristics.

\parahead{Communication Locality}
Baseline SFC-based placements (\S\ref{sec:amr:des:bg}) preserve spatial locality but hinder load
balance flexibility.
The actual runtime benefit of this locality had to be evaluated empirically, motivating exploring
locality-aware strategies.

We now describe the baseline infrastructure used and augmented by our policies in
\S\ref{sec:amr:des:bg}, followed by load-centric and locality-preserving policies in
\S\ref{sec:amr:des:lpt} and \S\ref{sec:amr:des:cdp}.
We conclude with CPLX (\S\ref{sec:amr:des:cplx}), a hybrid policy that enables tunable control over the
tradeoff between locality and load balance.
All policies were implemented in the Parthenon AMR framework~\cite{Grete2023:Parthenon} and are
directly available in Parthenon-based codes such as Phoebus~\cite{Barke2024:Phoebus} and
AthenaPK~\cite{Grete2023:Parthenon}.

\vspace{-5pt}

\subsection{Existing Placement Infrastructure}\label{sec:amr:des:bg}
% - Explain octree structure and why it's used - Show how SFCs preserve locality - Walk through
% baseline policy step by step - Highlight specific limitations that our policies address - Include
% a figure showing octree + SFC example

Modern block-based AMR codes use octrees and space-filling curves (SFCs) as foundational data
structures for mesh management.
This section describes these building blocks and the baseline placement policy, as understanding
their strengths and limitations is crucial for developing better policies.

\figoctree

\subsubsection{Octrees and Space-Filling Curves}

Octrees provide an efficient hierarchical representation of the mesh structure.
When a block needs to be refined, it is split into 8 ($2^3$) child blocks (in 3D).
Only leaf blocks of the octree participate in the simulation.
Octrees serve multiple functions in AMR codes---from mesh management to neighbor tracking to
block placement---with the latter enabled through easy construction of Z-order space-filling
curves via a depth-first traversal of the tree, as shown in \cref{fig:amr:bg:octree}.
Blocks are assigned sequential \emph{block IDs} in order of this traversal.
This ordering approximately preserves spatial locality, as blocks with nearby IDs are more likely
to be spatial neighbors.
Some locality, however, is inevitably lost as dimensionality reduction is inherently lossy.

\subsubsection{Baseline Placement Policy}

Placement computation is a component of the \emph{redistribution} process, invoked if the mesh
structure changes because of (de)refinement operations.
Redistribution proceeds in three steps: blocks are first assigned block IDs using Z-order SFCs, the
placement policy computes new block-rank mappings for all blocks, and finally blocks are migrated
to their new ranks using MPI P2P operations.

The baseline policy simply orders blocks by block ID and assigns contiguous ranges of $\lceil n/r
  \rceil$ or $\lfloor n/r \rfloor$ blocks to consecutive ranks, where $n$ is the total number of
blocks and $r$ is the number of ranks.
While better assignments are theoretically possible, this approach provides a reasonable balance
between compute load and communication locality by balancing block counts across ranks while
co-locating spatial neighbors.

\subsubsection{Augmenting Existing Infrastructure}

Octrees and space-filling curves are fundamental, well-studied components of modern AMR codes.
Rather than replace these proven building blocks, we augment them with mechanisms to handle
variable block costs and flexible allocation patterns.
While frameworks like Parthenon provide hooks for specifying per-block costs, these are typically
initialized to $1$ in practice---treating all blocks as computationally equal.
This simplified model persists partly because domain scientists rarely provide cost estimation
models for their simulations, and partly because the baseline policy's strict contiguous allocation
constraint limits its ability to accommodate variable costs.
Our changes to the infrastructure are minimal but enable significantly more
flexible placement policies:

\begin{enumerate}
  \item We populate the existing cost specification hooks with actual computation costs
        measured via telemetry
  \item We modify the block management infrastructure to support arbitrary
        (non-contiguous) block-to-rank mappings
\end{enumerate}

The second change required minor modifications to data structures but has no correctness
implications, as logical dependencies between blocks remain intact.
This allows for non-contiguous placements, but communication costs increase only if policies choose
to break locality.

\subsection{LPT: Prioritizing Load Balance}\label{sec:amr:des:lpt}
% - Brief overview of progression from LPT to CPLX

% Given reliable per-block costs and strict time constraints, we developed our placement policies
% through systematic evolution --- starting with simple, well-understood algorithms and carefully
% adding complexity only where beneficial. This section presents this evolution from the
% Longest-Processing-Time First (LPT) algorithm to our final hybrid policy.

% \subsubsection{LPT: Prioritizing Load Balance} - Intuition behind LPT - Algorithm description -
% 2-approximation guarantee (intuitive explanation) - Why it works surprisingly well in practice -
% Extension to multi-phase (brief, details in appendix) - Discussion of locality impact
We first investigated whether pure load balancing, ignoring communication locality entirely, could
provide meaningful runtime improvements.
The \emph{Longest-Processing-Time First} (LPT) algorithm offers a simple but powerful
approach---sort blocks by compute cost and assign each to the least loaded rank.
LPT is a classical greedy algorithm for makespan minimization with strong theoretical
guarantees---the makespan of an LPT solution is provably at most 4/3 times the optimal
makespan~\cite{Graha1969:LPT}.
In practice, LPT performs remarkably well; we could not obtain better solutions from a commercial
ILP solver~\cite{gurobi} despite letting it run for 200\,s.

LPT demonstrated surprisingly strong performance despite completely ignoring communication
locality.
While it did incur higher communication costs than the baseline approach, the improvements in load
balance more than compensated for this overhead.
This suggested that the value of locality preservation might be lower than initially assumed,
motivating us to explore this tradeoff further.

\vspace{-5pt}
\subsection{CDP: Preserving Locality}\label{sec:amr:des:cdp}
% - Motivation: cases where locality matters - Algorithm description with worked example - Analysis
% of tradeoffs vs LPT - Parallel implementation for scale - Why we need something better
After observing LPT perform well despite its locality costs, we wanted to explore the potential of
a locality-maximizing load-balanced placement.
To this end, we developed \emph{Contiguous-DP} (CDP), which uses dynamic programming to select
contiguous block ranges to improve load balance while preserving the locality patterns of the
baseline.

Consider a simulation with 10 blocks and 4 ranks.
With an average of 2.5 blocks per rank, CDP explores allocations like \texttt{[2,2,3,3]} or
\texttt{[2,3,2,3]}, finding one that minimizes makespan while maintaining contiguous allocations.
As a result, CDP has identical locality-preserving properties as baseline.

\parahead{Formal Description}
CDP solves the contiguous partitioning problem: given block costs $w_1, \dots, w_n$ (in SFC order),
partition them into $r$ contiguous segments to minimize the maximum segment sum (makespan).
Let $W[i] = \sum_{j=1}^{i} w_j$.
The minimum makespan $\text{DP}[i][k] $ for partitioning $i$ blocks among $k$ ranks follows
($\text{DP}[0][0] = 0$): \[ \text{DP}[i][k] = \min_{0 \le j < i} \max\left( \text{DP}[j][k-1] ,
  W[i] - W[j] \right) \]

The algorithm has complexity $O(n^2r)$ in general, where
$n$ is the number of blocks and $r$ is the number of ranks.
However, we optimize it to $O(nr)$ by considering only two contiguous chunk sizes: $\lceil n/r
  \rceil$ and $\lfloor n/r \rfloor$.
This optimization maintains solution quality while making CDP practical for AMR timescales.
This DP formulation guarantees an optimal makespan-minimizing placement within the explored chunk
sizes.

In our experiments, CDP improved upon baseline but only reached parity with LPT.
While CDP showed lower communication times, its higher synchronization times suggested that perfect
locality preservation might be unnecessarily restrictive, motivating our development of our hybrid
policy called CPLX, described next.

\parahead{Scaling CDP With Chunking}
The overhead of computing placements with CDP became noticeable at 4096 ranks, prompting us to
develop a parallel implementation using hierarchical chunking.
This approach divides blocks into $c$ contiguous chunks of approximately equal cost, then applies
CDP independently to each chunk using a subset of ranks.
At 4096 ranks with chunk size 512, this creates 8 parallel-processed chunks.

While chunking may not find the globally optimal CDP solution, this approximation has minimal
impact since CDP's output serves only as an intermediate step for CPLX.
The approach reduces placement overhead while retaining CDP's solution quality.

\subsection{CPLX: Combining the Best of Both}\label{sec:amr:des:cplx}

Our experiments with LPT and CDP revealed a clear tradeoff---LPT achieved better load balance but
higher communication costs, while CDP preserved locality at the expense of some load imbalance.
Neither policy consistently outperformed the other in end-to-end runtime, suggesting the potential
for a hybrid approach that could combine their strengths.

Developing such a hybrid proved challenging.
Our initial attempts to blend the policies produced unpredictable results---small sacrifices in
load balance did not translate to proportional gains in locality, and vice versa.
We eventually realized that it was easier to selectively break locality in a contiguous placement
than to restore locality in an arbitrary one.
This insight led to the CPLX design principle: start with a locality-preserving solution from CDP
and strategically apply LPT to address the most significant imbalances.

CPLX begins by computing an initial CDP placement that reuses the chunking mechanism for
scalability.
It then assesses the work allocated to different ranks by sorting them in descending order of load.
The policy selects X\% of ranks from either end of this sorted list---the most overloaded and
underloaded ranks---for rebalancing via LPT.
Including both ends is crucial as rebalancing needs both source and destination ranks to
effectively redistribute work.

The parameter X, ranging from 0 to 100\%, provides precise control over this locality disruption.
At X=0 (\texttt{CPL0}), no ranks participate in rebalancing, preserving the locality
characteristics of CDP.
At X=100 (\texttt{CPL100}), all ranks undergo LPT rebalancing, effectively becoming pure LPT.
Intermediate values create hybrid behaviors, only disrupting locality within the X\% most
imbalanced ranks while preserving it elsewhere.
This enables CPLX to systematically explore the complete locality-balance tradeoff space while
maintaining the performance requirements of AMR codes.
